\begin{pastexam}[label={px:gw3-096-ds-02}]{關務三等資訊處理 96 年 第二題}
（申論；20 分）

假設一個含有十個不相等鍵值（Key）的檔案，每個鍵都有對應使用頻率如下：
\begin{center}
\begin{tabular}{c|c}
鍵（Key） & 使用頻率（\%） \\ \hline
K$_{1}$ & 5 \\
K$_{2}$ & 10 \\
K$_{3}$ & 5 \\
K$_{4}$ & 20 \\
K$_{5}$ & 25 \\
K$_{6}$ & 5 \\
K$_{7}$ & 10 \\
K$_{8}$ & 5 \\
K$_{9}$ & 5 \\
K$_{10}$ & 10 \\
\end{tabular}
\end{center}
\begin{itemize}[leftmargin=2.2em, nosep]
\item[㈠] 若使用循序搜尋法（Sequential Search），試問搜尋一個鍵值（Key）的平均比較次數（Average Number of Comparisons）
為多少？（10 分）
\item[㈡] 若想有效減少平均比較次數，這些鍵值應重新安排（Arrangement），經過重新安排之後，試問其平均比較次數為多少？
（10 分）
\end{itemize}
\end{pastexam}

\begin{pastsoln}
答案：㈠ $5.3$ 次。 ㈡ 依使用頻率由高到低排列，$3.85$ 次。

\textbf{解讀。} 題目給的頻率合計 $100\%$，表示每次搜尋的鍵一定在檔案裡（成功搜尋）。 檔案依 K$_1$、K$_2$、…、K$_{10}$ 的
順序存放，循序搜尋從頭往後比，找到第 $i$ 個位置的鍵要比較 $i$ 次。 所以平均比較次數是
\[
\sum_{i=1}^{10} i \cdot p_i,
\]
$p_i$ 是放在第 $i$ 個位置的鍵被搜尋的機率（資結書第~1.8~節的平均情況分析）。

\textbf{㈠ 原來的順序。}
\[
\begin{aligned}
&1(0.05) + 2(0.10) + 3(0.05) + 4(0.20) + 5(0.25) \\
&\quad + 6(0.05) + 7(0.10) + 8(0.05) + 9(0.05) + 10(0.10) \\
&= 0.05 + 0.20 + 0.15 + 0.80 + 1.25 + 0.30 + 0.70 + 0.40 + 0.45 + 1.00 = 5.3.
\end{aligned}
\]

\textbf{㈡ 重新安排。} 把常被找的鍵放前面：依頻率由高到低排成
K$_5$ (25)、K$_4$ (20)、K$_2$、K$_7$、K$_{10}$ (各 10)、K$_1$、K$_3$、K$_6$、K$_8$、K$_9$ (各 5)，同頻率的次序不影響結果。
\[
\begin{aligned}
&1(0.25) + 2(0.20) + 3(0.10) + 4(0.10) + 5(0.10) + (6 + 7 + 8 + 9 + 10)(0.05) \\
&= 0.25 + 0.40 + 0.30 + 0.40 + 0.50 + 2.00 = 3.85.
\end{aligned}
\]

\textbf{這樣排是最好的。} 若某個排法有相鄰的兩個鍵，前面的頻率 $p$ 比後面的 $q$ 小，把兩者對調：前者多比一次、後者
少比一次，平均比較次數改變 $p - q < 0$，變小了。 所以最佳排法中頻率一定由高到低（交換論證，和資結書第~19.4~節
最短處理時間優先的證明相同）。 平均比較次數從 $5.3$ 降到 $3.85$。 頻率事先不知道時，可以用「找到就往前移」
（move-to-front）或「和前一個交換」（transpose）讓常用的鍵自己慢慢移到前面。
\end{pastsoln}
